% !TEX program = xelatex
% !TEX root = 笔记.tex
\documentclass[cn,hazy,blue,14pt,normal]{elegantnote}
\title{共形场论笔记}
\author{郑元昊}
\date{\zhtoday}
\institute{复旦大学物理系}

\input{packages.tex}

\begin{document}
\maketitle

\setcounter{tocdepth}{3} % 0: chapter/part, 1: section, 2: subsection, 3: subsubsection
\tableofcontents

\section{绪论}
\subsection{Why CFT?}
\begin{enumerate}
    \item CFT can emerge in "critical" phen: Ising model.
    \item Important to understand quantum gravity through string theory, AdS/CFT.
    \item CFT is truly formulated QFT non-perturbatively: Operator Algebra.
    \item CFT appears as fixed points of RG flow.
\end{enumerate}

\section{共形变换}
\subsection{时空的共形变换}
考虑一个$d$维时空流形$M$上的单参微分同胚群, 其对应的光滑矢量场为$\xi^a$, 时空度规$g_{ab}$(未指定号差, 可以Lorentz也可以Eucilidean)在其下的Lie导数为
\begin{align}
    \mathscr L_\xi g_{ab}=\xi^c\nabla_c g\omega_{ab}+g_{ac}\nabla_b\xi^c+g_{cb}\nabla_a\xi^c=\nabla_b\xi_a+\nabla_a\xi_b.
\end{align}

若保持时空度规不变, 即等距群, 则$\xi^a$满足
\begin{align}
    \nabla_{(a}\xi_{b)}=0.
\end{align}
即$\xi^a$为Killing矢量场.

在Minkovski时空中, 
\begin{align}
    g_{\mu\nu}=\rm{diag}(\pm1), 
\end{align}
容易解得Killing矢量场为
\begin{align}
    \xi^\mu=a^\mu+\omega^\mu_{~\nu}x^\nu,
\end{align}
其中
\begin{align}
    \omega_{\mu\nu}=\omega_{[\mu\nu]}.
\end{align}
这就是Poincare群, 总共有三部分组成: 平移、旋转、Boost.

如果我们放宽条件, 允许变换后的度规可以存在缩放, 即
\begin{align}
    g_{\mu\nu}'=g_{\mu\nu}-\mathscr L_\xi g_{\mu\nu}=\Omega g_{\mu\nu}=\exp{-2\sigma}g_{\mu\nu}\approx(1-2\sigma)g_{\mu\nu}
\end{align}
这就是共形变换. 不难发现, 从几何意义上说, 其实共性变换对称也就是保角变换.

若$\sigma$为常数$\lambda$, 则我们有
\begin{align}
    \partial_{(\mu}\xi_{\nu)}=\lambda g_{\mu\nu},
\end{align}
可以解得其解为通解+特解:
\begin{align}
    \xi^\mu=a^\mu+\omega^\mu_{~\nu}x^\nu+\lambda x^\mu.
\end{align}

如果再放宽条件, 我们让$\sigma=\sigma(x)$, 和坐标有关, 则
\begin{align}
    \partial_{(\mu}\xi_{\nu)}=\sigma g_{\mu\nu}\label{sp-c-eq}
\end{align}
将\expr{sp-c-eq}与$g_{\mu\nu}$缩并我们有
\begin{align}
    \partial_\mu\xi^\mu=d\sigma.\label{div-killing}
\end{align}
将\expr{sp-c-eq}与$\partial_\mu$缩并我们有
\begin{align}
    \Box\xi_\nu+\partial_\nu\partial_\mu\xi^\mu=\Box\xi_\nu+d\partial_\nu\sigma=2\partial_\nu\sigma,
\end{align}
即
\begin{align}
    \Box\xi_\nu=(2-d)\partial_\nu\sigma\label{sp-c-eq2}.
\end{align}
进一步将\expr{sp-c-eq2}与$\partial_\nu$缩并, 我们有
\begin{align}
    \Box(d\sigma)=(2-d)\Box\sigma
\end{align}
即
\begin{align}
    (1-d)\Box\sigma=0.
\end{align}

对于$d>1$, 我们有
\begin{align}
    \Box\sigma=0.
\end{align}

我们将$\partial_\mu$作用于\expr{sp-c-eq2}, 有
\begin{align}
    \Box\partial_\mu\xi_\nu=(2-d)\partial_\mu\partial_\nu\sigma.
\end{align}
同时全称化有, 
\begin{align}
    (2-d)\partial_\mu\partial_\nu\sigma=(2-d)\partial_{(\mu}\partial_{\nu)}\sigma=\Box\partial_{(\mu}\xi_{\nu)}=g_{\mu\nu}\Box\sigma=0.
\end{align}
即
\begin{align}
    (2-d)\partial_\mu\partial_\nu\sigma=0.
\end{align}

对于$d>2$的情况, 我们有
\begin{align}
    \partial_\mu\partial_\nu\sigma=0.
\end{align}
可以解得
\begin{align}
    \sigma=\lambda+2b_\mu x^\mu, 
\end{align}
其中$b_\mu$为某个常矢量. 然后代入求解$\xi^\mu$, 有
\begin{align}
    \partial_{(\mu}\xi_{\nu)}=(2b\cdot x+\lambda)g_{\mu\nu}.
\end{align}
不难得到其解
\begin{align}
    \xi^\mu=a^\mu+\omega^\mu_{~\nu}x^\nu+\lambda x^\mu+2(b\cdot x)x^\mu-x^2 b^\mu.
\end{align}
这就是共性变换对应的全部矢量场. 某一个特定矢量场定义了流形上的一个单参微分同胚群, 不难验证将这些单参微分同胚群取并集, 仍然具有封闭性、结合律、单位元以及逆元, 因此构成一个群, 即共性变换群.

其中, $a^\mu$代表平移(Translation), $\omega^\mu_{~\nu}$代表旋转(Rotation\&Boost), $\lambda$代表缩放(Dilation), $b^\mu$代表特殊共形变换(Special Conformal Transformation, SCT). 然后我们可以求这些无穷小变换对应的有限大坐标变换的表达式, 即微分方程
\begin{align}
    \de{x^\mu}{\tau}=\xi^\mu
\end{align}
在$\tau=1$时的解:
\begin{align}
    &\text{Translation:} x'^\mu=x^\mu+a^\mu,\\
    &\text{Boost:}  x'^\mu=\Lambda^\mu_{~\nu}x^\nu,\\
    &\text{Dilation:}  x'^\mu=\exp{\lambda}x^\mu,\\
    &\text{SCT:} x'^\mu=\frac{x^\mu-x^2b^\mu}{1-2b\cdot x+b^2 x^2}\label{sct-trans}.
\end{align}

其中\expr{sct-trans}的得到并不显然, 补充推导如下, 首先有微分方程, 
\begin{align}
    \de{x^\mu}{\tau}=2(b\cdot x)x^\mu-x^2b^\mu.
\end{align}
与$x^\mu$缩并我们有
\begin{align}
    \de{x^2}{\tau}=2(b\cdot x)x^2.
\end{align}
设
\begin{align}
    y^\mu=\frac{x^\mu}{x^2}, 
\end{align}
有
\begin{align}
    \de{y^\mu}{\tau}=-b^\mu
\end{align}
于是可得
\begin{align}
    \frac{x'^\mu}{x'^2}=\frac{x^\mu}{x^2}-b^\mu.
\end{align}
由此可见, SCT可以将其视为三个步骤: 以原点为中心内外翻转, 平移$-b^\mu$, 以原点为中心内外翻转. 最终化简我们有
\begin{align}
    x'^\mu=\frac{x^\mu-x^2b^\mu}{1-2b\cdot x+b^2 x^2}.
\end{align}
于是我们发现几个特殊点:
\begin{align}
    \begin{cases}
        & \text{origin}\to\text{origin}\\
        & \frac{b^\mu}{b^2}\to\infty\\
        & \infty\to-\frac{b^\mu}{b^2}.
    \end{cases}
\end{align}

\begin{theorem}\label{ct-3pt}
    设任意的六点$x_1, x_2, x_3$与$x_1', x_2', x_3'$, 存在共性变换将$x_1, x_2, x_3$变换为$x_1', x_2', x_3'$.
\end{theorem}
\begin{proof}
    首先利用平移, 可以将$x_1$平移到原点. 然后利用SCT, 可以将$x_2$变换到$\infty$, 这个时候$x_3$位于$x_3''$.
    
    类似地我们可以做到将$x_1'$平移到原点, $x_2'$变换到$\infty$, $x_3'$变换到$x_3'''$.

    如果我们对$x_3''$再做一个缩放与旋转, 一定可以将之变换到$x_3'''$. 由共性变换群的封闭性与逆元存在可知, 我们一定可以将$x_1, x_2, x_3$变换到$x_1', x_2', x_3'$.
\end{proof}

定理\ref{ct-3pt}告诉我们, 在CFT中计算关联函数时, 三点及以下的关联函数都是与点的位置关系没有任何关系的: 我们可以用共性变换将之变成任何形状!

\subsection{共形代数(Conformal Algebra)}
为了能够计算得到共形代数生成元之间的对易关系, 很自然的想法是我们寻找一个忠实表示, 在这个表示中计算对易子从而得到对易关系. 因此我们选取其对标量场的作用作为表示来考虑, 即我们有映射$G\to(\mathcal F\to\mathcal F)$: 设$g\in G, g=\exp{-i\omega_a K_a}$
\begin{align}
    g\to\phi_g, \phi_gf|_{p}\equiv \phi_{*t=1}f|_p.
\end{align}
其中, $\phi$对应的光滑矢量场即$\omega_a K_a$.

于是有
\begin{align}
    (1-i\epsilon a^\mu p_\mu-\frac i2\epsilon\omega^{\mu\nu}M_{\mu\nu}-i\epsilon\lambda D-i\epsilon b^\mu K_\mu)f(x^\mu)=f(x-\epsilon\xi^\mu).
\end{align}
所以可以计算得
\begin{align}
    & p_\mu=-i\partial_\mu, \\
    & M_{\mu\nu}=i(x_\mu\partial_\nu-x_\nu\partial_\mu), \\
    & D=-ix^\mu\partial_\mu, \\
    & K_\mu=i(x^2\partial_\mu-2x_\mu x^\nu\partial_\nu).
\end{align}

然后我们可以计算对易子
\begin{align}
    & [M^{\mu\nu}, M^{\rho\sigma}]=i(-g^{\mu\rho}M^{\nu\sigma}-g^{\sigma\nu}M^{\mu\rho}+g^{\mu\sigma}M^{\nu\rho}+g^{\rho\nu}M^{\mu\sigma}),\\
    & [M_{\mu\nu}, p_\rho]=-i(g_{\mu\rho}p_\nu-g_{\rho\nu}p_\mu),\\
    & [M_{\mu\nu}, K_\rho]=-i(g_{\mu\rho}K_\nu-g_{\rho\nu}K_\mu),\\
    & [D, p_\mu]=ip_\mu,\\
    & [D, K_\mu]=-iK_\mu,\\
    & [p_\mu, K_\nu]=-2i(M_{\mu\nu}+g_{\mu\nu}D),
\end{align}
这就是共形代数.

\begin{theorem}
    对于旋转对称群$\cong SO(m, n)$的流形, 其共形变换对称群$\cong SO(m+1, n+1)$.
\end{theorem}

\begin{proof}
    原本度规为$g_{\mu\nu}$, 设一个新的$m+n+2$维的度规
    \begin{align}
        \bar g_{MN}=g_{\mu\nu}\dx^\mu\dx^\nu+\dx^M\dx^M-\dx^N\dx^N,
    \end{align}
    其旋转对称群的Lie代数对易满足
    \begin{align}
        [L^{\mu\nu}, L^{\rho\sigma}]=i(-g^{\mu\rho}L^{\nu\sigma}-g^{\sigma\nu}L^{\mu\rho}+g^{\mu\sigma}L^{\nu\rho}+g^{\rho\nu}L^{\mu\sigma}).
    \end{align}
    如果令
    \begin{align}
        \begin{cases}
            D=L_{MN}=-L^{MN},\\
            p^\mu=L^{M\mu}+L^{N\mu},\\
            K^\mu=L^{M\mu}-L^{N\mu},\\
            M^{\mu\nu}=L^{\mu\nu}.
        \end{cases}
    \end{align}
    则可以使其满足全部的共形代数对易关系!
\end{proof}

\subsection{经典场的共形变换}
因为质量是和标度有关的, 因此只有无质量标量场具有标度不变性, 所以我们只考虑
\begin{align}
    \mathcal L=\frac12(\partial\phi)^2.
\end{align}

考虑无穷小微分同胚变换
\begin{align}
    x^\mu\to x'^\mu=x^\mu+\epsilon\xi^\mu.
\end{align}
我们可以将对标量场的变换总结为
\begin{align}
    \phi'(x')=\mathcal F\phi(x).
\end{align}
其中, $x\to x'$为时空的变换, 而$\mathcal F$则是标量场自身的变换. 于是可得
\begin{align}
    \phi(x)\to\phi'(x)=\phi(x)-\epsilon\xi^\mu\partial_\mu\phi+\epsilon\pa{\mathcal F}{\epsilon}(\phi).
\end{align}

那么剩下的问题就是我们该如何确定$\mathcal F(\phi)$? 我们考虑有限大标度变换$x\to\exp{-\lambda}x$, 我们想要物理规律与标度没关系, 则只能要求作用量$S$不变:
\begin{align}
    S'=\int\d^dx'\frac12(\partial'\phi')^2=S.
\end{align}
设$\phi'(x')=\exp{-\Delta\lambda}\phi(x)$, 根据
\begin{align}
    \d^dx'=\exp{d\lambda}\d^dx, \partial'=\exp{-\lambda}\partial
\end{align}
则我们有
\begin{align}
    S'=\exp{(-2\Delta-2+d)\lambda}S.
\end{align}
故
\begin{align}
    \Delta=\frac{d-2}{2}.
\end{align}

因此我们可得
\begin{align}
    \mathcal F(\phi)=\exp{\Delta\sigma}\phi\approx(1+\Delta\sigma)\phi, 
\end{align}
利用\expr{div-killing}我们有
\begin{align}
    \mathcal F(\phi)\approx\left(1+\frac{\Delta}d\partial_\mu\xi^\mu\right)\phi=\left(1+\frac{d-2}{2d}\partial_\mu\xi^\mu\right)\phi. 
\end{align}

因此对于无穷小共形变换, 我们有
\begin{align}
    & x^\mu\to x'^\mu=x^\mu+\epsilon\xi^\mu, \\
    & \phi(x)\to\phi'(x)=\phi(x)-\epsilon\xi^\mu\partial_\mu\phi-\epsilon\frac{d-2}{2d}\partial_\mu\xi^\mu\phi.
\end{align}

然后我们考虑这样变换作用量的变化
\begin{align}
    S'&=\int\d^dx'\mathcal L\left(\phi'(x'), \partial_\mu'\phi'(x')\right)\\
    &=\int\d^dx\left|\pa{x'}{x}\right|\mathcal L\left(\mathcal F\phi(x), \partial_\mu'\mathcal F\phi(x)\right).
\end{align}
利用
\begin{align}
    & |1+A|\approx 1+\Tr(A)\\
    & \mathcal F\phi(x)-\phi(x)\approx \epsilon\pa{\mathcal F}{\epsilon}\left(\phi(x)\right)
\end{align}
以及
\begin{align}
    \partial_\mu'(\mathcal F\phi)-\partial_\mu\phi&=(\delta_\mu^{~~\nu}-\partial_\mu\xi^\nu)\partial_\nu(\phi+\epsilon\pa{\mathcal F}{\epsilon}\left(\phi(x)\right))\\
    &=-\epsilon\partial_\mu\xi^\nu\partial_\nu\phi+\epsilon\partial_\mu\left(\pa{\mathcal F}\epsilon\right), 
\end{align}
我们有
\begin{align}
    S'&=\int\d^dx(1+\epsilon\partial_\mu\xi^\mu)\left[\mathcal L+\epsilon\pa{\mathcal L}\phi\pa{\mathcal F}\epsilon+\epsilon\pa{\mathcal L}{\partial_\mu\phi}\left(-\partial_\mu\xi^\nu\partial_\nu\phi+\partial_\mu\left(\pa{\mathcal F}{\epsilon}\right)\right)\right].
\end{align}
分部积分并代入EoM后, 我们有
\begin{align}
    \delta S=-\epsilon\int\d^dx T^{\mu\nu}\partial_{(\mu}\xi_{\nu)}.
\end{align}
其中$T^\mu\nu$为能动张量
\begin{align}
    T^{\mu\nu}=\partial^\mu\phi\partial^\nu\phi-\frac12(\partial\phi)^2g^{\mu\nu}.
\end{align}

\nocite{*}
\printbibliography[heading=bibintoc, title=\ebibname]

\end{document}